\uuid{JuCo}
\chapitre{Probabilité discrète}
\niveau{L2}
\module{Probabilité}
\sousChapitre{Probabilité et dénombrement}
\cpge{mpsi-mp2i.probabilites-univers-fini.c, pcsi.probabilites-variables-aleatoires.c, ptsi.probabilites-variables-aleatoires.c, mpsi-mp2i.probabilites-univers-fini.e, pcsi.probabilites-variables-aleatoires.e, ptsi.probabilites-variables-aleatoires.e, bcpst1.concepts-base-probabilites.b, bcpst1.variable-aleatoire-univers-fini.b}
\titre{Urnes, boules et indépendance}
\theme{probabilités}
\auteur{}
\datecreate{2023-01-24}
\organisation{AMSCC}
\difficulte{}
\contenu{

\texte{ 	On dispose de deux urnes. La première contient 3 boules noires et 7 boules blanches. La seconde contient 5 boules noires et 5 boules blanches. 

On choisit une urne au hasard et on tire successivement deux boules, avec remise, dans cette urne. On note les événements :
\begin{itemize}
	\item $B1$ : \og on obtient une boule blanche au premier tirage\fg{}
	\item $B2$ : \og on obtient une boule blanche au second tirage\fg{}
\end{itemize} }

\question{ Les événements $B1$ et $B2$ sont-ils indépendants ? }}
